\subsection{Entwurf einer Verstärkerschaltung mit unterschiedlichen Eingangsverstärkungen}
% Alt: Aufgabe 12
\Aufgabe{
  \label{Aufg12}
  Zwischen der Ausgangsspannung $U_\mathrm{a}$ und dem Eingangsspannungen $U_\mathrm{1}$ und $U_\mathrm{2}$ besteht der folgende
  funktionale Zusammenhang:
  $U_\mathrm{a} = 2 \ U_\mathrm{1} - 0,5 \ U_\mathrm{2}$
  Geben Sie eine mit Operationsverstärker(n) und Widerständen aufgebaute Schaltung zur Realisierung
  dieses funktionalen Zusammenhangs an. Die Operationsverstärker dürfen als ideale Operationsverstärker
  angenommen werden. Geben Sie Werte der verwendeten Widerstände an! Typische Widerstände für OPV-Schaltungen liegen im
  Bereich von $\mathrm{1\,k\Omega \ bis\ 100\,k\Omega}$.

}

\Loesung{
  \begin{figure}[H]
    \centering
    \input{Tikz/src/LsgOPV2Signale2.tex}\alttext{Die Abbildung zeigt eine Schaltung mit einem Operationsverstärker, der mit einem Widerstandsnetzwerk am invertierenden und am nichtinvertierenden Eingang beschaltet ist. Der Operationsverstärker besitzt einen invertierenden Eingang (−), einen nichtinvertierenden Eingang (+) sowie einen Ausgang.
      Der Ausgang des Operationsverstärkers befindet sich auf der rechten Seite. Die Ausgangsspannung \(U_\mathrm{a}\) wird gegen Masse gemessen. Am invertierenden Eingang liegt das Eingangssignal \(U_{\mathrm{e1}}\) an. Dieses wird über den Widerstand \(R_1\) dem invertierenden Eingang zugeführt. Zusätzlich ist der invertierende Eingang über den Widerstand \(R_3\) mit dem Ausgang verbunden. Dadurch entsteht eine Gegenkopplung vom Ausgang zum invertierenden Eingang.
      Der nichtinvertierende Eingang ist ebenfalls über ein Widerstandsnetzwerk beschaltet. Das Eingangssignal \(U_{\mathrm{e2}}\) wird über den Widerstand \(R_2\) zugeführt. Zudem ist der nichtinvertierende Eingang über den Widerstand \(R_4\) mit Masse verbunden.
      Die beiden Eingangsspannungen \(U_{\mathrm{e1}}\) und \(U_{\mathrm{e2}}\) werden jeweils gegen Masse gemessen. Der gemeinsame Massebezug ist unten dargestellt.}
    \label{LsgOPV2Signale2}
  \end{figure}

  Der funktionale Zusammenhang zwischen der Ausgangsspannung \( U_a \) und den Eingangsspannungen \( U_{e1} \) und \( U_{e2} \) lautet:
  \[
    U_\mathrm{a} = 2 U_\mathrm{e1} - 0{,}5 U_\mathrm{e2}
  \]

  \textbf{1. Allgemeine Formel für die Ausgangsspannung eines Subtrahierers:}
  \[
    U_\mathrm{a} = U_\mathrm{e2} \cdot \frac{R_1 + R_3}{R_2 + R_4} \cdot \frac{R_4}{R_1} - U_\mathrm{e1} \cdot \frac{R_3}{R_1}
  \]
  Vergleich mit der Vorgabe \( U_\mathrm{a} = 2 U_\mathrm{e1} - 0{,}5 U_\mathrm{e2} \) ergibt folgende Bedingungen: \\
  \\
  - Verstärkungsfaktor für \( U_\mathrm{e1} \): \( \frac{R_3}{R_1} = 0{,}5 \) \\
  \\
  - Verstärkungsfaktor für \( U_\mathrm{e2} \): \( \frac{R_1 + R_3}{R_2 + R_4} \cdot \frac{R_4}{R_1} = 2 \)


  \textbf{2. Vorgehen:} \\
  - Setze \( R_3 = 1\,\mathrm{k}\Omega \) (typischer Wert). \\
  - Berechne die restlichen Widerstände.


  \textbf{3. Berechnung der Widerstände:}

  Berechnung von \( R_1 \) (aus \( \frac{R_3}{R_1} = 0{,}5 \)):
  \[
    \frac{R_3}{R_1} = 0{,}5 \quad \Rightarrow \quad R_1 = \frac{R_3}{0{,}5} = \frac{1\,\mathrm{k}\Omega}{0{,}5} = 2\,\mathrm{k}\Omega
  \]

  Berechnung von \( R_2 \) und \( R_4 \) (aus \( \frac{R_1 + R_3}{R_2 + R_4} \cdot \frac{R_4}{R_1} = 2 \)):
  Setze \( R_4 = 1\,\mathrm{k}\Omega \):
  \[
    \frac{R_1 + R_3}{R_2 + R_4} \cdot \frac{R_4}{R_1} = 2 \quad \Rightarrow \quad \frac{3\,\mathrm{k}\Omega}{R_2 + 1\,\mathrm{k}\Omega} \cdot \frac{1\,\mathrm{k}\Omega}{2\,\mathrm{k}\Omega} = 2
  \]
  Multipliziere beide Seiten:
  \[
    \frac{3}{2} = 2 \cdot (R_2 + 1\,\mathrm{k}\Omega)
  \]
  Umstellen nach \( R_2 \):
  \[
    R_2^2 + R_2 - \frac{3}{2} = 0
  \]


  \textbf{4. Lösung mit der PQ-Formel:}

  Die quadratische Gleichung hat die Form:
  \[
    R_2^2 + p \cdot R_2 + q = 0 \quad \text{mit} \quad p = 1, \, q = -\frac{3}{2}
  \]
  PQ-Formel:
  \[
    R_2 = -\frac{p}{2} \pm \sqrt{\left( \frac{p}{2} \right)^2 - q}
  \]
  Einsetzen der Werte:
  \[
    R_2 = -\frac{1}{2} \pm \sqrt{\left( \frac{1}{2} \right)^2 - \left( -\frac{3}{2} \right)} = -\frac{1}{2} \pm \sqrt{\frac{1}{4} + \frac{3}{2}}
  \]
  \[
    R_2 = -\frac{1}{2} \pm \sqrt{\frac{7}{4}} = -\frac{1}{2} \pm \frac{\sqrt{7}}{2}
  \]
  Da \( R_2 \) positiv sein muss:
  \[
    R_2 = -\frac{1}{2} + \frac{\sqrt{7}}{2} \approx 1{,}82\,\mathrm{k}\Omega
  \]


  \textbf{5. Ergebnis der Widerstände:}
  \begin{align*}
    R_1 & = 2\,\mathrm{k}\Omega      \\
    R_2 & = 1{,}82\,\mathrm{k}\Omega \\
    R_3 & = 1\,\mathrm{k}\Omega      \\
    R_4 & = 1\,\mathrm{k}\Omega
  \end{align*}
  Siehe: Abschnitt \ref{fig: Verschaltung Verstaerker} und Tabelle \ref{tab:Grundschaltungen}
}